\section{Manipulador Robótico}
Neste trabalho, o subsistema do manipulador robótico (MR) é acionado na fase final da missão de RVD/B.
O MR está rigidamente acoplado à base do chaser, possuindo  configuração 4R1P.
Em ambiente de microgravidade, o movimento do MR não sofre ação de peso, porém, por causa da conservação do momento angular, acaba produzindo reações na base, que alteram a atitude e o tensor de inércia do conjunto.

% Foi adotada a modelagem clássica de robótica com elos rígidos e juntas sem perdas ou associação com hardware direto, além disso definiu-se os referenciais de acordo com a convenção DH (clássica).
% O estudo do MR abrange:
% (i) cinemática (direta e diferencial),
% (ii) dinâmica e acoplamento base--manipulador,
% (iii) controle em juntas e por fim, 
% (iv) simulação numérica com restrições e limites.

\subsection{Convenção}
Neste trabalho, o termo \emph{ferramenta}, \emph{tool} e \emph{efetuador} são equivalentes e designam o mesmo elemento físico na extremidade ativa do MR.
O seu referencial foi denominado $\{E\}$ e se encontra no ponto central da ferramenta para expressar a posição e a orientação da ponta.
Portanto, ${}^{0}T_{E}$ denota a pose de
$\{E\}$ em relação à base $\{0\}$.
Além disso, adotou-se $\{0\}\equiv\{base\}$ (sem offset global) e,
quando aplicável, inclui-se a transformação da ferramenta:
${}^5 T_{tool} = {}^0 T_{5}\,{}^0 T_{tool}$.

\input{2 Formulacao do problema/2.8 Manipulador robotico/2.8.1 Modelagem cinematica.tex}
\subsection{Cinemática Diferencial e Jacobiano}

A cinemática diferencial descreve como as velocidades de junta geram as velocidades da ferramenta em um referencial escolhido.
Foi utilizada a matriz Jacobiana $\mathbf{J}(\vec q)$ para mapear
$\dot{\vec q}$ em velocidade linear e angular da ferramenta.
% Este mapeamento é a base para:
% (i) controle no espaço de tarefa,
% (ii) análise de singularidades e
% (iii) conversão entre forças cartesianas e torques de junta.

\subsubsection{Velocidades cartesianas e Jacobiana}
A partir do referencial da base $\{0\}$, a velocidade (linear e angular) cartesiana da ferramenta é obtida diretamente a partir das velocidades das juntas pela Jacobiana geométrica.
Para o MR configuração 4R1P, no mesmo sistema de coordenadas, as contribuições de cada junta ao movimento da ferramenta:
\begin{equation}
\label{eq:xdot_J_qdot}
{}^{0}\!\dot{\vec x}_{E} \;=\;
\begin{bmatrix}
{}^{0}\!\vec v_{E} \\
{}^{0}\!\vec \omega_{E}
\end{bmatrix}
= \mathbf{J}(\vec q)\,\dot{\vec q},
\end{equation}
em que $\dot{\vec q}$ reúne as velocidades.

\subsubsection{Construção do Jacobiano}
As colunas do Jacobiano são construídas no referencial $\{0\}$ a partir dos eixos
${}^{0}\!\vec z_{\,i-1}$ e
dos vetores posição ${}^{0}\!\vec p_{\,\cdot}$ obtidos pela cinemática direta.

Para a junta revoluta:
\begin{equation}
\begin{cases}
J_{v,i} =
{}^{0}\!\vec{z}_{\,i-1} \times 
\left({}^{0}\!\vec{p}_{\,E} - {}^{0}\!\vec{p}_{\,i-1}\right), \\
J_{\omega,i} =
{}^{0}\!\vec{z}_{\,i-1}.
\end{cases}
\end{equation}

Para a junta prismática:
\begin{equation}
\begin{cases}
J_{v,i} = {}^{0}\!\vec{z}_{\,i-1}, \\
J_{\omega,i} = \vec 0.
\end{cases}
\end{equation}

Portanto, cada coluna de $\mathbf{J}$ representa a contribuição de cada junta para o movimento linear e angular do MR.

\subsubsection{Efeito da Base}
Quando a base está fixa, usa-se diretamente
\eqref{eq:xdot_J_qdot}.
Para base flutuante, a velocidade da ferramenta é a soma do efeito do movimento da base com o efeito do movimento das juntas:
\begin{equation}
{}^{0}\!\dot{\vec x}_{E} =
\mathbf{J}_{base}(\cdot)\,{}^{0}\!\dot{\vec x}_{B}
\;+\;
\mathbf{J}(\vec q)\,\dot{\vec q}.
\end{equation}
Neste trabalho, $\mathbf{J}_{base}(\cdot)$ é tratada de forma agregada na seção de acoplamento base--manipulador; aqui, assume-se base fixa no uso de $\mathbf{J}(\vec q)$ para controle.

\subsubsection{Mapeamento de torques e forças}
O mapeamento entre as forças--momentos (\emph{wrench}) aplicados na ferramenta e os torques nas juntas é dado por:

\begin{equation}
\vec\tau = \mathbf{J}^T (\vec q)\,\vec w_{E},
\end{equation}

em que a \emph{wrench} na ferramenta é
\begin{equation}
\vec w_{E} =
\begin{bmatrix}
\vec f_{E} \\
\vec m_{E}
\end{bmatrix},
\end{equation}
com $\vec f_{E}$ (N) e $\vec m_{E}$ (N$\cdot$m) expressos no mesmo referencial e aplicados no mesmo ponto $E$.
Portanto, ($\mathbf{J}^{\top}\vec w_{E}$) transforma as forças--momentos cartesianos desejados (ou medidos) na ferramenta em torques equivalentes nas juntas.

\subsubsection{Singularidades}
Uma singularidade ocorre quando pequenas variações desejadas no efetuador $\{E\}$ exigem velocidades de junta excessivas em alguma direção.
Isso acontece quando algum termo de menor valor de $\mathbf{J}(\vec q)$ tende a zero e o número de condição cresce.
A pseudo-inversa amortecida pode ser utilizada, e é expressa por:

\begin{equation}
\label{eq:pseudoinversa}
\dot{\vec q} =
\mathbf{J}^T\,\big(\mathbf{J}\,\mathbf{J}^T
+ \lambda^2 I\big)^{-1}\,{}^{0}\!\dot{\vec x}_{E,d},
\end{equation}
com $\lambda>0$.
% Para o MR deste trabalho, de configuração 4R1P, a dimensão de tarefa é $5$ (posição 3D e dois ângulos relevantes à captura), o que reduz o risco de singularidades.
\input{2 Formulacao do problema/2.8 Manipulador robotico/2.8.3 Modelagem dinamica.tex}
\subsection{Acoplamento base--manipulador}
\label{sec:acoplamento_base_manipulador}

As variáveis são para o acoplamento dos movimentos entre a base e o MR:
\begin{equation}
\label{eq:nu_def}
\nu
\;=\;
\begin{bmatrix}
{}^{0}\!\vec v_{B} \\[2pt]
{}^{0}\!\vec \omega_{B} \\[2pt]
\dot{\vec q}
\end{bmatrix},
\end{equation}
onde ${}^{0}\!\vec v_{B}$ e
${}^{0}\!\vec \omega_{B}$
são as velocidades linear e angular da base expressas em $\{0\}$, e
$\dot{\vec q}$ são as velocidades das juntas.

\subsection{Jacobiano composto da base e do MR}
A velocidade cartesiana da ferramenta no mesmo sistema de coordenadas $\{0\}$ passa a ter duas contribuições, sendo elas:

\begin{equation}
\label{eq:xdot_composto}
{}^{0}\!\dot{\vec x}_{E}
\;=\;
\underbrace{\mathbf{J}_{\!B}(\cdot)}_{\text{efeito da base}}\,
\begin{bmatrix}
{}^{0}\!\vec v_{B} \\
{}^{0}\!\vec \omega_{B}
\end{bmatrix}
\;+\;
\underbrace{\mathbf{J}(\vec q)}_{\text{efeito das juntas}}\,\dot{\vec q},
\end{equation}

onde $\mathbf{J}_{\!B}(\cdot)$ depende apenas da posição relativa entre a base $\{B\}$ e
o efetuador $\{E\}$, enquanto que
$\mathbf{J}(\vec q)$ é o Jacobiano geométrico do manipulador.

\subsection{Matriz de inércia e o acoplamento}
Considerando como base livre, as equações de movimento podem ser descritas por:

\begin{equation}
\label{eq:EOM_particionada_compacta}
\mathbf M_{\text{tot}}(\vec q\, )\,\dot{\nu}
\;+\;
\mathbf C_{\text{tot}}(\vec q,\nu)\,\nu
\;=\;
\vec\tau_{\text{tot}}
\;+\;
\vec w_{\text{ext}}.
\end{equation}

Detalhando cada termo, obtém-se:
\begin{equation}
\label{eq:Mtot_def}
\mathbf M_{\text{tot}}(\vec q)
=
\begin{bmatrix}
\mathbf M_{BB}(\vec q) & \mathbf M_{BM}(\vec q) \\
\mathbf M_{MB}(\vec q) & \mathbf M_{MM}(\vec q)
\end{bmatrix}
\end{equation}

\begin{equation}
\label{eq:Ctot_def}
\mathbf C_{\text{tot}}(\vec q,\nu)
=
\begin{bmatrix}
\mathbf C_{BB}(\cdot) & \mathbf C_{BM}(\cdot) \\
\mathbf C_{MB}(\cdot) & \mathbf C_{MM}(\cdot)
\end{bmatrix}
\end{equation}

% Vetores empilhados (opcional, ajuda a reduzir largura na EOM)
\begin{equation}
\label{eq:stacked_defs1}
\dot{\nu} =
\begin{bmatrix}
{}^{0} \dot{\vec v}_{B} \\
{}^{0} \dot{\vec \omega}_{B} \\
\ddot{\vec q}
\end{bmatrix},
\qquad
\nu =
\begin{bmatrix}
{}^{0} \vec v_{B} \\
{}^{0} \vec \omega_{B} \\
\dot{\vec q}
\end{bmatrix},
\end{equation}

% Vetores empilhados (opcional, ajuda a reduzir largura na EOM)
\begin{equation}
\label{eq:stacked_tau}
\vec\tau_{\text{tot}} =
\begin{bmatrix}
\vec \tau_{B} \\
\vec \tau
\end{bmatrix},
\end{equation}

\begin{equation}
\label{eq:stacked_wrench}
\vec w_{\text{ext}} =
\begin{bmatrix}
{}^{\text{ext}}\vec w_{B} \\
{}^{\text{ext}}\vec w_{E}
\end{bmatrix}.
\end{equation}

A equação \eqref{eq:EOM_particionada_compacta} descreve a dinâmica acoplada base--manipulador.
A matriz de inércia \eqref{eq:Mtot_def} agrega a inércia translacional e rotacional da base, a inércia no espaço das juntas e os termos de acoplamento dinâmico entre base e braço; além disso, é simétrica e definida positiva.

A matriz \eqref{eq:Ctot_def} reúne os termos de Coriolis/centrífugos do sistema base e manipulador, incluindo efeitos cruzados.
% Deve satisfazer a propriedade estrutural $\dot{\mathbf M}_{\text{tot}}-2\,\mathbf C_{\text{tot}}$ anti-simétrica, assegurando coerência energética.

Os vetores $\nu$ e $\dot{\nu}$ definidos em \eqref{eq:stacked_defs1} são, respectivamente, as velocidades linear e angular da base, além das velocidades articulares; e as correspondentes acelerações.

Os vetores em \eqref{eq:stacked_tau} são as entradas de atuação (esforços na base e torques/força nas juntas).
Por sua vez, \eqref{eq:stacked_wrench} explicita as excitações externas (força--momento na base e na ferramenta), todas expressas no mesmo referencial e referidas aos pontos $B$ e $E$, respectivamente.

\subsubsection{Momento angular}
Para descrever a interação do momento angular do sistema em torno do centro de massa, utiliza-se:

\begin{equation}
\label{eq:centroidal}
\vec h_{c}
\;=\;
\mathbf A_{c}(\vec q)\,\nu,
\qquad
\dot{\vec h}_{c} \;=\; \vec w^{\text{ext}}_{c}.
\end{equation}

Na ausência de momentos externos
(${}^{\text{ext}}\vec w_{c}=\vec 0$),
o momento angular no centro de massa é conservado: movimentos das juntas alteram $\vec h_{c}$ e, por reação, geram velocidade/atitude na base.
A matriz $\mathbf A_{c}(\vec q)$ é obtida por composição rígida (corpos articulados) e depende apenas da configuração.